% =========================================================
% เฉลยแบบฝึกหัดบทที่ 2 ระบบควบคุม
% หมายเหตุ: ไฟล์นี้เป็น LaTeX snippet สำหรับนำไป \input{} หรือคัดลอกเพิ่มในเอกสารเดิม
% ไม่รวม \documentclass, \usepackage และ \begin{document}
% =========================================================

\section{เฉลยแบบฝึกหัดบทที่ 2 ระบบควบคุม}

\subsection{ส่วนที่ 1: สมการเชิงอนุพันธ์ (ODE)}

\subsubsection{ข้อ 1: สมการเชิงอนุพันธ์ของกระแสในวงจร $RL$}

กำหนดให้
\[
R=5\ \Omega,\qquad L=2\ \mathrm{H},\qquad v(t)=10u(t)\ \mathrm{V}
\]

สำหรับวงจร $RL$ อนุกรม ใช้กฎแรงดันของเคอร์ชอฟฟ์
\[
v(t)=v_R(t)+v_L(t)
\]
โดย
\[
v_R(t)=Ri(t)
\]
และ
\[
v_L(t)=L\frac{di(t)}{dt}
\]
ดังนั้น
\[
L\frac{di(t)}{dt}+Ri(t)=v(t)
\]
แทนค่า $L=2$, $R=5$ และ $v(t)=10u(t)$ จะได้
\[
2\frac{di(t)}{dt}+5i(t)=10u(t)
\]
ดังนั้นสมการเชิงอนุพันธ์คือ
\[
\boxed{2\frac{di(t)}{dt}+5i(t)=10u(t)}
\]

ถ้ากำหนดให้กระแสเริ่มต้นเป็นศูนย์ คือ $i(0)=0$ จะได้รูปมาตรฐาน
\[
\frac{di(t)}{dt}+\frac{5}{2}i(t)=5u(t)
\]
สำหรับ $t\ge 0$ จะมี $u(t)=1$ ดังนั้น
\[
\frac{di(t)}{dt}+\frac{5}{2}i(t)=5
\]
ค่ากระแสสุดท้ายคือ
\[
i(\infty)=\frac{v}{R}=\frac{10}{5}=2\ \mathrm{A}
\]
ค่าคงเวลาของวงจรคือ
\[
\tau=\frac{L}{R}=\frac{2}{5}=0.4\ \mathrm{s}
\]
ดังนั้นผลตอบสนองของกระแสคือ
\[
\boxed{i(t)=2\left(1-e^{-2.5t}\right)u(t)}
\]

\subsubsection{ข้อ 2: ค่าคงเวลาและอัตราขยายของระบบอันดับหนึ่ง}

กำหนดสมการ
\[
2\frac{dy}{dt}+6y=12u(t)
\]
จัดรูปให้อยู่ในรูปมาตรฐานของระบบอันดับหนึ่ง โดยหารทั้งสมการด้วย $6$
\[
\frac{1}{3}\frac{dy}{dt}+y=2u(t)
\]
เปรียบเทียบกับรูปมาตรฐาน
\[
\tau\frac{dy}{dt}+y=Ku(t)
\]
จะได้
\[
\boxed{\tau=\frac{1}{3}\ \mathrm{s}}
\]
และ
\[
\boxed{K=2}
\]

\subsection{ส่วนที่ 2: การแปลงลาปลาซและการแปลงลาปลาซผกผัน}

\subsubsection{ข้อ 3: การแปลงลาปลาซของฟังก์ชันเวลา}

ต้องการหา
\[
\mathcal{L}\left\{5e^{-3t}+2\cos(4t)-3\sin(2t)\right\}
\]
ใช้สูตรพื้นฐาน
\[
\mathcal{L}\{e^{-at}\}=\frac{1}{s+a}
\]
\[
\mathcal{L}\{\cos \omega t\}=\frac{s}{s^2+\omega^2}
\]
\[
\mathcal{L}\{\sin \omega t\}=\frac{\omega}{s^2+\omega^2}
\]
ดังนั้น
\[
\mathcal{L}\{5e^{-3t}\}=\frac{5}{s+3}
\]
\[
\mathcal{L}\{2\cos(4t)\}=\frac{2s}{s^2+16}
\]
และ
\[
\mathcal{L}\{-3\sin(2t)\}=-\frac{6}{s^2+4}
\]
จึงได้
\[
\boxed{
\mathcal{L}\left\{5e^{-3t}+2\cos(4t)-3\sin(2t)\right\}
=
\frac{5}{s+3}+\frac{2s}{s^2+16}-\frac{6}{s^2+4}
}
\]

\subsubsection{ข้อ 4: การแปลงลาปลาซผกผันโดยแยกเศษส่วนย่อย}

ต้องการหา
\[
\mathcal{L}^{-1}\left\{\frac{3s+7}{(s+2)(s+5)}\right\}
\]
แยกเศษส่วนย่อยเป็น
\[
\frac{3s+7}{(s+2)(s+5)}=\frac{A}{s+2}+\frac{B}{s+5}
\]
คูณทั้งสองข้างด้วย $(s+2)(s+5)$
\[
3s+7=A(s+5)+B(s+2)
\]
จัดรูปเป็น
\[
3s+7=(A+B)s+(5A+2B)
\]
เทียบสัมประสิทธิ์
\[
A+B=3
\]
\[
5A+2B=7
\]
จาก $B=3-A$ แทนในสมการที่สอง
\[
5A+2(3-A)=7
\]
\[
3A=1
\]
\[
A=\frac{1}{3}
\]
ดังนั้น
\[
B=3-\frac{1}{3}=\frac{8}{3}
\]
จึงได้
\[
\frac{3s+7}{(s+2)(s+5)}=\frac{1/3}{s+2}+\frac{8/3}{s+5}
\]
แปลงกลับได้
\[
\boxed{y(t)=\frac{1}{3}e^{-2t}+\frac{8}{3}e^{-5t}}
\]

\subsubsection{ข้อ 5: การแปลงลาปลาซผกผันกรณีมีขั้วซ้ำ}

ต้องการหา
\[
\mathcal{L}^{-1}\left\{\frac{s+4}{(s+1)(s+3)^2}\right\}
\]
แยกเศษส่วนย่อยเป็น
\[
\frac{s+4}{(s+1)(s+3)^2}=\frac{A}{s+1}+\frac{B}{s+3}+\frac{C}{(s+3)^2}
\]
คูณทั้งสมการด้วย $(s+1)(s+3)^2$
\[
s+4=A(s+3)^2+B(s+1)(s+3)+C(s+1)
\]
แทน $s=-1$
\[
3=4A
\]
ดังนั้น
\[
A=\frac{3}{4}
\]
แทน $s=-3$
\[
1=-2C
\]
ดังนั้น
\[
C=-\frac{1}{2}
\]
จากสัมประสิทธิ์ของ $s^2$ ได้
\[
A+B=0
\]
จึงได้
\[
B=-\frac{3}{4}
\]
ดังนั้น
\[
\frac{s+4}{(s+1)(s+3)^2}=\frac{3/4}{s+1}-\frac{3/4}{s+3}-\frac{1/2}{(s+3)^2}
\]
แปลงกลับได้
\[
\boxed{y(t)=\frac{3}{4}e^{-t}-\frac{3}{4}e^{-3t}-\frac{1}{2}te^{-3t}}
\]

\subsubsection{ข้อ 6: การแปลงลาปลาซผกผันโดยวิธี Complete the Square}

ต้องการหา
\[
\mathcal{L}^{-1}\left\{\frac{2s+5}{s^2+6s+13}\right\}
\]
จัดรูปตัวส่วนโดยทำกำลังสองสมบูรณ์
\[
s^2+6s+13=(s+3)^2+4
\]
จัดตัวเศษให้อยู่ในรูปของ $s+3$
\[
2s+5=2(s+3)-1
\]
ดังนั้น
\[
\frac{2s+5}{s^2+6s+13}
=
\frac{2(s+3)-1}{(s+3)^2+4}
\]
แยกเป็น
\[
\frac{2(s+3)}{(s+3)^2+4}-\frac{1}{(s+3)^2+4}
\]
เนื่องจาก $4=2^2$ จะได้
\[
\frac{2(s+3)}{(s+3)^2+2^2}-\frac{1}{(s+3)^2+2^2}
\]
พจน์ที่สองเขียนเป็น
\[
\frac{1}{(s+3)^2+2^2}=\frac{1}{2}\frac{2}{(s+3)^2+2^2}
\]
ดังนั้น
\[
\boxed{y(t)=2e^{-3t}\cos(2t)-\frac{1}{2}e^{-3t}\sin(2t)}
\]

\subsection{ส่วนที่ 3: IVT, FVT และการแยกเศษส่วนย่อย}

\subsubsection{ข้อ 7: ระบบที่มีฟังก์ชันถ่ายโอนและอินพุตแบบ Unit Step}

กำหนดให้
\[
G(s)=\frac{10}{(s+2)(s+5)}
\]
และอินพุตเป็น Unit Step
\[
R(s)=\frac{1}{s}
\]

\paragraph{(ก) หา $Y(s)$}

ใช้ความสัมพันธ์
\[
Y(s)=G(s)R(s)
\]
ดังนั้น
\[
Y(s)=\frac{10}{(s+2)(s+5)}\cdot\frac{1}{s}
\]
จึงได้
\[
\boxed{Y(s)=\frac{10}{s(s+2)(s+5)}}
\]

\paragraph{(ข) หา $y(\infty)$ ด้วย Final Value Theorem}

ใช้ทฤษฎีบทค่าสุดท้าย
\[
y(\infty)=\lim_{s\to 0}sY(s)
\]
แทนค่า
\[
y(\infty)=\lim_{s\to 0}s\frac{10}{s(s+2)(s+5)}
\]
ตัด $s$ ได้
\[
y(\infty)=\lim_{s\to 0}\frac{10}{(s+2)(s+5)}
\]
แทน $s=0$
\[
y(\infty)=\frac{10}{(2)(5)}=1
\]
ดังนั้น
\[
\boxed{y(\infty)=1}
\]

\paragraph{(ค) ตรวจสอบเงื่อนไขการใช้ FVT}

ต้องพิจารณาขั้วของ $sY(s)$
\[
sY(s)=\frac{10}{(s+2)(s+5)}
\]
ขั้วคือ
\[
s=-2,\qquad s=-5
\]
ขั้วทั้งสองอยู่ในครึ่งซ้ายของระนาบ $s$ จึงสามารถใช้ FVT ได้
\[
\boxed{\text{สามารถใช้ FVT ได้ เพราะขั้วของ }sY(s)\text{ อยู่ในครึ่งซ้ายทั้งหมด}}
\]

เพื่อยืนยันด้วยการแปลงกลับ พิจารณา
\[
Y(s)=\frac{10}{s(s+2)(s+5)}
\]
แยกเศษส่วนย่อยเป็น
\[
Y(s)=\frac{A}{s}+\frac{B}{s+2}+\frac{C}{s+5}
\]
หา $A$
\[
A=\left.\frac{10}{(s+2)(s+5)}\right|_{s=0}=1
\]
หา $B$
\[
B=\left.\frac{10}{s(s+5)}\right|_{s=-2}=-\frac{5}{3}
\]
หา $C$
\[
C=\left.\frac{10}{s(s+2)}\right|_{s=-5}=\frac{2}{3}
\]
ดังนั้น
\[
Y(s)=\frac{1}{s}-\frac{5}{3}\frac{1}{s+2}+\frac{2}{3}\frac{1}{s+5}
\]
แปลงกลับได้
\[
\boxed{y(t)=1-\frac{5}{3}e^{-2t}+\frac{2}{3}e^{-5t}}
\]
เมื่อ $t\to\infty$ จะได้ $y(\infty)=1$ ตรงกับผลจาก FVT

\subsubsection{ข้อ 8: การหา $y(0^+)$ ด้วย IVT และ $y(\infty)$ ด้วย FVT}

กำหนดให้
\[
\frac{Y(s)}{R(s)}=\frac{6}{s^2+5s+6}
\]
และ
\[
R(s)=\frac{1}{s}
\]
ดังนั้น
\[
Y(s)=\frac{6}{s^2+5s+6}\cdot\frac{1}{s}
\]
เนื่องจาก
\[
s^2+5s+6=(s+2)(s+3)
\]
จึงได้
\[
\boxed{Y(s)=\frac{6}{s(s+2)(s+3)}}
\]

\paragraph{หา $y(0^+)$ ด้วย Initial Value Theorem}

ใช้ทฤษฎีบทค่าเริ่มต้น
\[
y(0^+)=\lim_{s\to\infty}sY(s)
\]
ดังนั้น
\[
y(0^+)=\lim_{s\to\infty}s\frac{6}{s(s+2)(s+3)}
\]
ตัด $s$ ได้
\[
y(0^+)=\lim_{s\to\infty}\frac{6}{(s+2)(s+3)}=0
\]
ดังนั้น
\[
\boxed{y(0^+)=0}
\]

\paragraph{หา $y(\infty)$ ด้วย Final Value Theorem}

ใช้
\[
y(\infty)=\lim_{s\to 0}sY(s)
\]
จะได้
\[
y(\infty)=\lim_{s\to 0}s\frac{6}{s(s+2)(s+3)}
\]
ตัด $s$ ได้
\[
y(\infty)=\lim_{s\to 0}\frac{6}{(s+2)(s+3)}
\]
แทน $s=0$
\[
y(\infty)=\frac{6}{(2)(3)}=1
\]
ดังนั้น
\[
\boxed{y(\infty)=1}
\]

\paragraph{ยืนยันด้วยการแปลงกลับ}

แยกเศษส่วนย่อย
\[
Y(s)=\frac{6}{s(s+2)(s+3)}=\frac{A}{s}+\frac{B}{s+2}+\frac{C}{s+3}
\]
หา $A$
\[
A=\left.\frac{6}{(s+2)(s+3)}\right|_{s=0}=1
\]
หา $B$
\[
B=\left.\frac{6}{s(s+3)}\right|_{s=-2}=-3
\]
หา $C$
\[
C=\left.\frac{6}{s(s+2)}\right|_{s=-3}=2
\]
ดังนั้น
\[
Y(s)=\frac{1}{s}-\frac{3}{s+2}+\frac{2}{s+3}
\]
แปลงกลับได้
\[
\boxed{y(t)=1-3e^{-2t}+2e^{-3t}}
\]
ตรวจสอบค่าเริ่มต้น
\[
y(0^+)=1-3+2=0
\]
ตรวจสอบค่าสุดท้าย
\[
y(\infty)=1
\]

\subsection{ส่วนที่ 4: การแก้สมการเชิงอนุพันธ์ด้วยการแปลงลาปลาซ}

\subsubsection{ข้อ 9: การแก้สมการอันดับหนึ่งด้วยการแปลงลาปลาซ}

กำหนดสมการ
\[
\frac{dy}{dt}+5y=10u(t)
\]
โดยมีเงื่อนไขเริ่มต้น
\[
y(0)=2
\]
ใช้ลาปลาซทั้งสองข้าง
\[
\mathcal{L}\left\{\frac{dy}{dt}\right\}+5\mathcal{L}\{y(t)\}=\mathcal{L}\{10u(t)\}
\]
โดย
\[
\mathcal{L}\left\{\frac{dy}{dt}\right\}=sY(s)-y(0)
\]
และ
\[
\mathcal{L}\{10u(t)\}=\frac{10}{s}
\]
แทนค่า $y(0)=2$
\[
sY(s)-2+5Y(s)=\frac{10}{s}
\]
จัดรูป
\[
(s+5)Y(s)=\frac{10}{s}+2
\]
เขียนเป็นเศษส่วนเดียวกัน
\[
(s+5)Y(s)=\frac{10+2s}{s}
\]
ดังนั้น
\[
Y(s)=\frac{2s+10}{s(s+5)}
\]
ดึงตัวประกอบร่วม
\[
Y(s)=\frac{2(s+5)}{s(s+5)}
\]
ตัด $(s+5)$ ได้
\[
Y(s)=\frac{2}{s}
\]
แปลงกลับได้
\[
\boxed{y(t)=2u(t)}
\]
หรือสำหรับ $t\ge 0$
\[
\boxed{y(t)=2}
\]

เนื่องจากค่าเริ่มต้น $y(0)=2$ เท่ากับค่าคงตัวสุดท้ายของระบบ จึงไม่มีพจน์ทรานเซียนต์เกิดขึ้น กราฟของ $y(t)$ เป็นเส้นตรงแนวนอนที่ $y=2$ ตั้งแต่ $t=0$ เป็นต้นไป

\subsubsection{ข้อ 10: ระบบอันดับสองและประเภทการหน่วง}

กำหนดสมการ
\[
\frac{d^2y}{dt^2}+4\frac{dy}{dt}+4y=8u(t)
\]
โดยมี
\[
y(0)=0,\qquad y'(0)=0
\]

\paragraph{การหาประเภทการหน่วง}

พิจารณาตัวส่วนของระบบ
\[
s^2+4s+4
\]
เปรียบเทียบกับรูปมาตรฐานของระบบอันดับสอง
\[
s^2+2\zeta\omega_n s+\omega_n^2
\]
เทียบสัมประสิทธิ์
\[
\omega_n^2=4
\]
จึงได้
\[
\omega_n=2
\]
และ
\[
2\zeta\omega_n=4
\]
แทน $\omega_n=2$
\[
2\zeta(2)=4
\]
\[
\zeta=1
\]
ดังนั้นระบบเป็นระบบหน่วงวิกฤต
\[
\boxed{\zeta=1}
\]
\[
\boxed{\text{ระบบเป็นแบบ Critically Damped System}}
\]

\paragraph{การแก้สมการด้วยลาปลาซ}

ใช้ลาปลาซทั้งสองข้าง
\[
\mathcal{L}\{y''\}+4\mathcal{L}\{y'\}+4\mathcal{L}\{y\}=\mathcal{L}\{8u(t)\}
\]
โดย
\[
\mathcal{L}\{y''\}=s^2Y(s)-sy(0)-y'(0)
\]
และ
\[
\mathcal{L}\{y'\}=sY(s)-y(0)
\]
แทนค่าเริ่มต้น $y(0)=0$ และ $y'(0)=0$
\[
s^2Y(s)+4sY(s)+4Y(s)=\frac{8}{s}
\]
จัดรูป
\[
(s^2+4s+4)Y(s)=\frac{8}{s}
\]
เนื่องจาก
\[
s^2+4s+4=(s+2)^2
\]
ดังนั้น
\[
Y(s)=\frac{8}{s(s+2)^2}
\]

แยกเศษส่วนย่อย
\[
\frac{8}{s(s+2)^2}=\frac{A}{s}+\frac{B}{s+2}+\frac{C}{(s+2)^2}
\]
คูณด้วย $s(s+2)^2$
\[
8=A(s+2)^2+Bs(s+2)+Cs
\]
แทน $s=0$
\[
8=4A
\]
จึงได้
\[
A=2
\]
แทน $s=-2$
\[
8=-2C
\]
จึงได้
\[
C=-4
\]
จากสัมประสิทธิ์ของ $s^2$ ได้
\[
A+B=0
\]
ดังนั้น
\[
B=-2
\]
จึงได้
\[
Y(s)=\frac{2}{s}-\frac{2}{s+2}-\frac{4}{(s+2)^2}
\]
แปลงกลับได้
\[
\boxed{y(t)=2-2e^{-2t}-4te^{-2t}}
\]
หรือจัดรูปเป็น
\[
\boxed{y(t)=2\left[1-(1+2t)e^{-2t}\right]}
\]

\subsubsection{ข้อ 11: วงจร $RLC$ อนุกรม}

กำหนดให้
\[
R=2\ \Omega,\qquad L=1\ \mathrm{H},\qquad C=\frac{1}{5}\ \mathrm{F}
\]
ต่อกับแรงดัน
\[
v(t)=5u(t)
\]

\paragraph{(ก) สมการเชิงอนุพันธ์ของ $v_C(t)$}

สำหรับวงจร $RLC$ อนุกรม
\[
v(t)=v_L(t)+v_R(t)+v_C(t)
\]
โดย
\[
i(t)=C\frac{dv_C(t)}{dt}
\]
ดังนั้น
\[
v_R(t)=Ri(t)=RC\frac{dv_C(t)}{dt}
\]
และ
\[
v_L(t)=L\frac{di(t)}{dt}=LC\frac{d^2v_C(t)}{dt^2}
\]
จึงได้สมการ
\[
LC\frac{d^2v_C(t)}{dt^2}+RC\frac{dv_C(t)}{dt}+v_C(t)=v(t)
\]
แทนค่า $L=1$, $R=2$, $C=\frac{1}{5}$
\[
\frac{1}{5}\frac{d^2v_C}{dt^2}+\frac{2}{5}\frac{dv_C}{dt}+v_C=5u(t)
\]
คูณทั้งสมการด้วย $5$
\[
\boxed{\frac{d^2v_C}{dt^2}+2\frac{dv_C}{dt}+5v_C=25u(t)}
\]

\paragraph{(ข) แปลงลาปลาซ}

สมมติสภาวะเริ่มต้นเป็นศูนย์
\[
v_C(0)=0,\qquad \frac{dv_C(0)}{dt}=0
\]
ใช้ลาปลาซกับสมการ
\[
\frac{d^2v_C}{dt^2}+2\frac{dv_C}{dt}+5v_C=25u(t)
\]
จะได้
\[
s^2V_C(s)+2sV_C(s)+5V_C(s)=\frac{25}{s}
\]
ดังนั้น
\[
(s^2+2s+5)V_C(s)=\frac{25}{s}
\]
จึงได้
\[
\boxed{V_C(s)=\frac{25}{s(s^2+2s+5)}}
\]
เนื่องจาก
\[
s^2+2s+5=(s+1)^2+4
\]
จึงเขียนได้เป็น
\[
\boxed{V_C(s)=\frac{25}{s\left[(s+1)^2+4\right]}}
\]

\paragraph{(ค) หา $v_C(t)$}

แยกเศษส่วนย่อย
\[
\frac{25}{s(s^2+2s+5)}=\frac{A}{s}+\frac{Bs+C}{s^2+2s+5}
\]
คูณด้วย $s(s^2+2s+5)$
\[
25=A(s^2+2s+5)+s(Bs+C)
\]
จัดรูป
\[
25=(A+B)s^2+(2A+C)s+5A
\]
เทียบสัมประสิทธิ์
\[
5A=25
\]
จึงได้
\[
A=5
\]
จาก
\[
A+B=0
\]
ได้
\[
B=-5
\]
และจาก
\[
2A+C=0
\]
ได้
\[
C=-10
\]
ดังนั้น
\[
V_C(s)=\frac{5}{s}+\frac{-5s-10}{s^2+2s+5}
\]
จัดตัวส่วนและตัวเศษใหม่
\[
s^2+2s+5=(s+1)^2+4
\]
\[
-5s-10=-5(s+1)-5
\]
ดังนั้น
\[
V_C(s)=\frac{5}{s}-\frac{5(s+1)}{(s+1)^2+4}-\frac{5}{(s+1)^2+4}
\]
เนื่องจาก $4=2^2$ และ
\[
\frac{5}{(s+1)^2+4}=\frac{5}{2}\frac{2}{(s+1)^2+2^2}
\]
จึงได้
\[
V_C(s)=\frac{5}{s}-\frac{5(s+1)}{(s+1)^2+2^2}-\frac{5}{2}\frac{2}{(s+1)^2+2^2}
\]
แปลงกลับได้
\[
\boxed{v_C(t)=5-5e^{-t}\cos(2t)-\frac{5}{2}e^{-t}\sin(2t)}
\]
หรือ
\[
\boxed{v_C(t)=5-e^{-t}\left[5\cos(2t)+\frac{5}{2}\sin(2t)\right]}
\]

\paragraph{(ง) หา $v_C(\infty)$ ด้วย FVT}

ใช้ Final Value Theorem
\[
v_C(\infty)=\lim_{s\to 0}sV_C(s)
\]
จาก
\[
V_C(s)=\frac{25}{s(s^2+2s+5)}
\]
ดังนั้น
\[
v_C(\infty)=\lim_{s\to 0}s\frac{25}{s(s^2+2s+5)}
\]
ตัด $s$ ได้
\[
v_C(\infty)=\lim_{s\to 0}\frac{25}{s^2+2s+5}
\]
แทน $s=0$
\[
v_C(\infty)=\frac{25}{5}=5
\]
ดังนั้น
\[
\boxed{v_C(\infty)=5\ \mathrm{V}}
\]

\section{เฉลยคำถามทบทวน}

\subsection{คำถามที่ 1: เหตุผลที่ใช้การแปลงลาปลาซในการแก้สมการเชิงอนุพันธ์}

การแปลงลาปลาซมีประโยชน์มากในระบบควบคุม เพราะช่วยเปลี่ยนปัญหาจากโดเมนเวลาไปเป็นโดเมน $s$ ทำให้การวิเคราะห์ง่ายขึ้น ข้อได้เปรียบสำคัญมีดังนี้

\begin{enumerate}
    \item \textbf{เปลี่ยนสมการเชิงอนุพันธ์เป็นสมการพีชคณิต}
    
    สมการเชิงอนุพันธ์ เช่น
    \[
    \frac{dy}{dt}+ay=bu(t)
    \]
    เมื่อแปลงลาปลาซจะกลายเป็นสมการพีชคณิตในรูป
    \[
    (s+a)Y(s)=\text{input term}
    \]
    ทำให้แก้ได้ง่ายกว่า เพราะไม่ต้องใช้วิธีหาผลเฉลยทั่วไปและผลเฉลยเฉพาะแบบคลาสสิก

    \item \textbf{รวมเงื่อนไขเริ่มต้นได้โดยตรง}
    
    การแปลงลาปลาซสามารถนำค่าเริ่มต้น เช่น $y(0)$ และ $y'(0)$ เข้าไปในสมการได้ทันที เช่น
    \[
    \mathcal{L}\left\{\frac{dy}{dt}\right\}=sY(s)-y(0)
    \]
    จึงเหมาะกับระบบจริงที่มักมีพลังงานสะสมเริ่มต้น เช่น แรงดันเริ่มต้นของตัวเก็บประจุ หรือกระแสเริ่มต้นของตัวเหนี่ยวนำ

    \item \textbf{เหมาะสำหรับการวิเคราะห์ระบบควบคุม}
    
    ในระบบควบคุมนิยมใช้ฟังก์ชันถ่ายโอน
    \[
    G(s)=\frac{Y(s)}{R(s)}
    \]
    ซึ่งอยู่ในโดเมน $s$ การใช้ลาปลาซทำให้สามารถวิเคราะห์เสถียรภาพ ตำแหน่งขั้ว ผลตอบสนองต่อ Step Input และค่าคงตัวสุดท้ายได้ง่าย

    \item \textbf{ใช้กับอินพุตมาตรฐานได้สะดวก}
    
    อินพุตที่ใช้บ่อย เช่น $u(t)$, $\delta(t)$, $\sin \omega t$ และ $e^{-at}$ มีตารางลาปลาซรองรับอยู่แล้ว ทำให้หาผลตอบสนองของระบบได้รวดเร็ว
\end{enumerate}

\subsection{คำถามที่ 2: Region of Convergence (ROC)}

Region of Convergence หรือ ROC คือบริเวณของค่า $s$ ในระนาบเชิงซ้อนที่ทำให้อินทิกรัลลาปลาซลู่เข้า การแปลงลาปลาซนิยามโดย
\[
F(s)=\int_{0}^{\infty}f(t)e^{-st}\,dt
\]
โดยทั่วไป
\[
s=\sigma+j\omega
\]
ดังนั้น ROC คือช่วงของ $\sigma$ ที่ทำให้อินทิกรัลมีค่าจำกัด หรือกล่าวได้ว่า
\[
\boxed{\text{ROC คือบริเวณบนระนาบ }s\text{ ที่ทำให้ Laplace Transform มีค่าอยู่จริง}}
\]

ในระบบควบคุมโดยทั่วไป มักใช้ลาปลาซด้านเดียว หรือ One-sided Laplace Transform
\[
\mathcal{L}\{f(t)\}=\int_{0^-}^{\infty}f(t)e^{-st}\,dt
\]
ซึ่งเหมาะกับระบบ causal หรือระบบที่เริ่มพิจารณาที่ $t=0$ ดังนั้นจึงมักให้ความสำคัญกับตำแหน่งขั้วของระบบมากกว่า ROC โดยตรง เพราะตำแหน่งขั้วบอกเสถียรภาพของระบบได้ เช่น ถ้าขั้วทุกตัวอยู่ซีกซ้ายของระนาบ $s$ ระบบจะมีแนวโน้มเสถียร

\[
\boxed{\text{ในงานระบบควบคุมมักพิจารณา Poles และ Stability มากกว่า ROC โดยตรง}}
\]

\subsection{คำถามที่ 3: เงื่อนไขการใช้ Final Value Theorem}

Final Value Theorem คือ
\[
\lim_{t\to\infty}y(t)=\lim_{s\to 0}sY(s)
\]
หรือ
\[
\boxed{y(\infty)=\lim_{s\to 0}sY(s)}
\]

FVT ใช้ได้เมื่อขั้วของ $sY(s)$ ต้องอยู่ในครึ่งซ้ายของระนาบ $s$ ทั้งหมด กล่าวคือ ผลตอบสนองของระบบต้องเข้าสู่ค่าคงตัวจริงเมื่อเวลามากขึ้น

ตัวอย่างกรณีที่ใช้ได้ เช่น
\[
Y(s)=\frac{1}{s(s+2)}
\]
ดังนั้น
\[
sY(s)=\frac{1}{s+2}
\]
มีขั้วที่ $s=-2$ ซึ่งอยู่ซีกซ้าย จึงใช้ FVT ได้
\[
y(\infty)=\lim_{s\to 0}\frac{1}{s+2}=\frac{1}{2}
\]

ตัวอย่างกรณีที่ FVT ให้คำตอบผิด พิจารณา
\[
Y(s)=\frac{1}{s^2+1}
\]
ซึ่งเป็นลาปลาซของ
\[
y(t)=\sin t
\]
ถ้าใช้ FVT จะได้
\[
y(\infty)=\lim_{s\to 0}sY(s)=\lim_{s\to 0}\frac{s}{s^2+1}=0
\]
แต่ในความจริง $\sin t$ ไม่ลู่เข้าสู่ค่าคงตัว เพราะแกว่งตลอดเวลา ดังนั้น FVT ใช้ไม่ได้ เนื่องจากขั้วอยู่ที่
\[
s=\pm j
\]
ซึ่งอยู่บนแกนจินตภาพ

\[
\boxed{\text{FVT ใช้ไม่ได้กับระบบที่มีการแกว่งไม่จางหาย}}
\]

\subsection{คำถามที่ 4: เปรียบเทียบการแก้สมการเชิงอนุพันธ์ในโดเมนเวลากับโดเมน $s$}

\begin{center}
\begin{tabular}{p{0.45\textwidth}p{0.45\textwidth}}
\hline
\textbf{การแก้ในโดเมนเวลา} & \textbf{การแก้ในโดเมน $s$} \\
\hline
ทำงานกับสมการเชิงอนุพันธ์โดยตรง เช่น $\frac{dy}{dt}+ay=bu(t)$ & แปลงสมการเชิงอนุพันธ์เป็นสมการพีชคณิตในรูปของ $Y(s)$ \\
ต้องหาผลเฉลยทั่วไปและผลเฉลยเฉพาะ & ใช้การจัดรูปพีชคณิตและแยกเศษส่วนย่อย \\
เงื่อนไขเริ่มต้นต้องนำไปใช้หลังจากได้ผลเฉลยทั่วไป & เงื่อนไขเริ่มต้นถูกรวมเข้าไปในสมการลาปลาซโดยตรง \\
เหมาะกับสมการง่าย ๆ และต้องการเข้าใจพฤติกรรมในเวลาโดยตรง & เหมาะกับระบบควบคุม วงจร และระบบที่มีฟังก์ชันถ่ายโอน \\
วิเคราะห์เสถียรภาพอาจทำได้ยากกว่า & วิเคราะห์เสถียรภาพได้จากตำแหน่งขั้วของระบบ \\
ผลลัพธ์อยู่ในรูป $y(t)$ โดยตรง & ได้ $Y(s)$ ก่อน แล้วจึงแปลงกลับเป็น $y(t)$ \\
ไม่สะดวกเมื่อต้องเจอกับ Step, Impulse หรือระบบซับซ้อน & สะดวกมากกับอินพุตมาตรฐาน เช่น Step, Impulse และ Sinusoidal \\
\hline
\end{tabular}
\end{center}

สรุปคือ การแก้ในโดเมนเวลาเหมาะกับการมองพฤติกรรมโดยตรง ส่วนการแก้ในโดเมน $s$ เหมาะกับการวิเคราะห์และออกแบบระบบควบคุม

\subsection{คำถามที่ 5: ระบบอันดับสองที่มีขั้ว $s=-2\pm3j$}

ระบบอันดับสองมาตรฐานมีขั้วอยู่ที่
\[
s=-\zeta\omega_n\pm j\omega_n\sqrt{1-\zeta^2}
\]
โจทย์ให้ขั้วเป็น
\[
s=-2\pm 3j
\]
ดังนั้นส่วนจริงคือ
\[
-\zeta\omega_n=-2
\]
จึงได้
\[
\zeta\omega_n=2
\]
ส่วนจินตภาพคือ
\[
\omega_n\sqrt{1-\zeta^2}=3
\]
อีกวิธีหนึ่งคือใช้ความสัมพันธ์จากขั้วเชิงซ้อน
\[
\omega_n=\sqrt{\sigma^2+\omega_d^2}
\]
โดย
\[
\sigma=2,\qquad \omega_d=3
\]
ดังนั้น
\[
\omega_n=\sqrt{2^2+3^2}
\]
\[
\omega_n=\sqrt{13}
\]
จึงได้
\[
\boxed{\omega_n=\sqrt{13}\ \mathrm{rad/s}}
\]
หรือประมาณ
\[
\omega_n\approx 3.606\ \mathrm{rad/s}
\]
จาก
\[
\zeta\omega_n=2
\]
จะได้
\[
\zeta=\frac{2}{\omega_n}=\frac{2}{\sqrt{13}}
\]
ดังนั้น
\[
\boxed{\zeta=\frac{2}{\sqrt{13}}}
\]
หรือประมาณ
\[
\zeta\approx 0.555
\]
เนื่องจาก
\[
0<\zeta<1
\]
ระบบจึงเป็นระบบหน่วงต่ำกว่าวิกฤต หรือ Underdamped System
\[
\boxed{\text{ระบบนี้เป็น Underdamped System}}
\]
